Un café está a $72\,{}^\circ\mathrm C$ en aire a $20\,{}^\circ\mathrm C$. Tras $2$ minutos está a $67.6\,{}^\circ\mathrm C$. Use la ley de Newton para inferir la constante de enfriamiento y predecir la temperatura a los $7$ minutos. Compare la predicción con la medida $61.5\,{}^\circ\mathrm C$.

Se supone temperatura uniforme en el café, ambiente constante, constante $k$ positiva e invariable y ausencia de otras aportaciones de calor. Los dos datos usados para inferir $k$ se tratan como exactos dentro del cálculo del modelo.
